\section{Typed cosets and flag transitivity}
The following construction is Tits' coset construction, recalled as
Theorem~3.1 of Jacobs--Leemans. No finiteness of the group is needed here.
Let $I$ be a finite type set, $G$ a group, and $(H_i)_{i\in I}$ subgroups.
Take the \emph{disjoint typed union} of the right-coset sets
$H_i\backslash G$; the type remains part of an object even when two
subgroups happen to be equal. Declare two typed cosets incident when
their underlying subsets of $G$ intersect.

\begin{proposition}\label{tits}
This is a reflexive symmetric incidence structure, with no incident
distinct objects of the same type. The cosets $(H_i)_{i\in I}$ form a
chamber. Right multiplication by $G$ preserves incidence and types and
is transitive on flags of each prescribed type of size at most two.
It need not be transitive on larger flags, nor need every flag extend
to a chamber without an additional hypothesis.
\end{proposition}
\begin{proof}
Right cosets of a fixed subgroup are either equal or disjoint. This
gives the same-type assertion; reflexivity and symmetry follow from
intersection. The displayed cosets all contain the identity and include
one object of every type. Multiplication on the right is a bijection
of $G$ preserving intersections. Every object is a translate of its
standard coset. If $H_i g$ and $H_j h$ are incident, choose a common
element $a$; the two cosets equal $H_i a,H_j a$, respectively. Thus this
pair is the translate by $a$ of the standard pair. The empty flag is
fixed. These arguments prove exactly the asserted transitivity.
\end{proof}

In the finite algorithms below, full flag transitivity is proved by
checking all eligible-extension equalities, not inferred from this
pair-transitivity statement. A nontrivial kernel of the coset action
is allowed. Inversion identifies this right-coset convention with the
left-coset convention used in the forest application, so no extra
geometric premise is introduced when changing conventions.
